\section{Discussion}
\label{sec:discussion}

\subsection{Cross-View Aggregation of Program Reasoning}
Forward and backward tasks provide complementary views of the same program, gold value, and grader. Discrepancies between these views are highly informative: a method that improves forward performance while degrading backward performance demonstrates task-specific overfitting rather than generalized program comprehension. Evaluating either view in isolation obscures this phenomenon—as shown in Table~\ref{tab:main}, all adaptations appear beneficial in the forward direction. True progress requires cross-view aggregation.

We propose a strict aggregation criterion: true improvements in comprehension must meet or exceed the untuned model's performance across \emph{all} views and models. In our study, only the sign-consensus merge (Section~\ref{sec:results}) satisfies this across all eight models, despite never being the top performer in any single view. This criterion is suitably conservative (filtering out spurious improvements) and computationally free, requiring only bidirectional evaluation on existing behavior-prediction benchmarks.

\subsection{Choosing an Adaptation Strategy}
No single adaptation strategy dominates across all conditions. For fixed, known deployment distributions, the anchored objective and mixture models are optimal, yielding genuine forward-pass gains of 3--6 points. However, under distributional shift (e.g., updated obfuscators, unseen transform compositions, or novel query formats), the sign-consensus merge proves superior. It is the only method that never significantly underperforms the untuned baseline in any configuration, requiring only a single adapter and no additional inference overhead. While the mixture model achieves comparable out-of-distribution (OOD) performance, it requires eight adapters and an inference gate---and the learned gate offers no empirical advantage over a uniform prior on held-out families.

\subsection{The Limits of Training Breadth}
Intuitively, pooled training across multiple transforms should yield generalizable models. Our results contradict this in both OOD regimes. On the held-out family, pooled training underperforms all modular approaches composed from the same data across all eight models, and even underperforms a baseline adapter trained solely on clean code in six of them. In the backward direction, it degrades below the untuned model on four of the six readable models. Rather than learning generalized invariants, pooled training appears to memorize surface-level transformations—features that unseen obfuscators readily modify.

While training breadth remains effective for in-distribution data (performing within a few points of the optimal method on seen compositions), its benefits do not generalize. Crucially, the identical training dataset avoids this confinement when learned as separate specialists and composed post-hoc. The critical factor is not data volume or model capacity, but parameter isolation: learning conditions in distinct weight spaces versus a single shared representation. In static environments, this distinction is negligible; in dynamic or adversarial environments, it is paramount.